% Part: lambda-calculus
% Chapter: introduction
% Section: minimization

\documentclass[../../../include/open-logic-section]{subfiles}

\begin{document}

\olfileid[fr]{lam}{int}{min}
\olsection{Les fonctions \texorpdfstring{$\lambd$}{λ}-définissables et la minimisation}

\begin{lem}
Supposons que $f(x,y)$ soit totale et $\lambd$-définissable.
Soit $g$ la fonction donnée par
\[
g(x)\simeq\umin{y}{f(x,y)}.
\]
Alors $g$ est $\lambd$-définissable.
\end{lem}

\begin{proof}
Voici l'idée. Étant donné~$x$, nous utiliserons
un combinateur de point fixe~$Y$ pour définir
une fonction $h_x(n)$ qui cherche un~$y$ à partir
de~$n$; alors $g(x)=h_x(0)$. La recherche
séquentielle voulue s'exprime par
\[
h_x(n)\simeq
\begin{cases}
n & \text{si $f(x,n)=0$},\\
h_x(n+1) & \text{si $f(x,n)$ est défini et non nul},\\
\text{indéfini} & \text{si $f(x,n)$ est indéfini}.
\end{cases}
\]

Passons aux détails. Puisque $f$ est
$\lambd$-définissable par hypothèse, un
terme~$F$ la définit. Rappelons le terme~$D$ :
$D(M,N,\num{0})\red M$ et
$D(M,N,\num{1})\red N$. La même définition
donne $D(M,N,\num{r})\red N$ pour tout~$r\ge1$.
Choisissons des variables $x$, $h$ et~$z$ distinctes, hors de $\FV{F}$. Pour définir $h_x$,
cherchons un terme~$H$, qui dépend de~$x$, tel que
\[
H(\num{n})\equiv
D(\num{n},H(\fn{Succ}(\num{n})),F(x,\num{n})).
\]
Prenons pour~$Y$ le combinateur de Turing
de la section précédente. Définissons d'abord
\[
U\ident\lambd[h][\lambd[z][D(z,h(\fn{Succ}(z)),F(x,z))]],
\]
puis $H\ident YU$. Si~$F$ est fermé, la seule
variable libre de~$H$ est~$x$. Comme
le combinateur de Turing vérifie
$Y U\red U(Y U)$, nous avons
\[
H=YU\red U(YU)=U(H).
\]
En particulier, pour tout entier naturel~$n$,
\begin{align*}
H(\num{n}) & \red U(H,\num{n}) \\
& \red D(\num{n},H(\fn{Succ}(\num{n})),F(x,\num{n})),
\end{align*}
comme voulu. Si nous remplaçons~$x$ par un
numéral~$\num{m}$, cette expression se réduit
à~$\num{n}$ lorsque $F(\num{m},\num{n})$
se réduit à~$\num{0}$; elle se réduit à
$H(\fn{Succ}(\num{n}))$ lorsque
$F(\num{m},\num{n})$ se réduit à un autre
numéral.

Enfin, posons $G\ident\lambd[x][H(\num{0})]$.
Si la recherche trouve une première valeur nulle, $G(\num{m})$
se réduit à $\num{g(m)}$. S'il n'y en a aucune, la réduction de
tête enchaîne les tests sans atteindre une forme normale de tête.
Pour justifier ce dernier point, on applique l'argument du
\olref[lam][rep][par]{thm:partial-recursive-lambda} à la même
sélection $D$ sur des numéraux issus d'un test total.
\end{proof}

Avec un test partiel, l'expression $D(M,N,F(x,\num{n}))$ peut
sélectionner une branche alors que le troisième argument n'a pas
de forme normale. Posons
$B\ident\lambd[p][\lambd[q][p\Omega]]$,
où $\Omega\ident(\lambd[z][zz])(\lambd[z][zz])$,
et $F(x,n)\ident D(B,\num{0},n)$. Alors $F(x,\num{0})$
n'a pas de forme normale et $F(x,\num{n})\red\num{0}$ pour
tout~$n\ge1$. Pourtant, $D(\num{0},N,B)\red N$, car
$K(N)$ efface $\Omega$. La recherche passe alors au test suivant,
alors que la minimisation stricte est indéfinie dès ce premier
test. La fermeture pour la minimisation stricte découle de
la caractérisation par les fonctions récursives partielles.

\end{document}
